\cotefile{157-2}
\scanpage{307}

\mspage{48}
% archive sheet number: 306

6. \uline{\unsure{Compositions|Composition} des} $\underline{\mathrm{Ty}}$

On a une loi de composition (partout définie) dans $\operatorname{Ob}(\underline{\mathrm{Ty}})=\mathrm{Ty}$, c'est celle du monoïde libre à deux générateurs $\Delta_+,\Delta_-$. Cette loi de composition s'interprète, au niveau \deleted{des intervalles} catégorique (et non plus \unsure{seulement} combinatoire) comme une somme amalgamée
\[
\begin{aligned}
&\deleted{I_\tau I_{\tau'}=\unsure{I_{\tau,\tau'}|I_{\tau\tau'}}}\\
&I_{\tau\tau'}=I_\tau\amalg_{(e,\underline{1}_\tau,\underline{0}_{\tau'})}I_{\tau'}
\end{aligned}
\]
où $e$ est la catégorie unité, et où $\underline{1}_\tau,\underline{0}_{\tau'}$ représentent \unsure{maintenant} les morphismes $e\unsure{\text{ --- }|\longrightarrow}I_\tau$, $e\longrightarrow I_{\tau'}$ ayant comme image l'élément but et l'élément source respectivement de l'intervalle orienté envisagé. On a \unsure{donc} un carré cocartésien dans $(\mathrm{Cat})$
\diagram{\begin{tikzcd}[ampersand replacement=\&,column sep=small,row sep=small]
\& e \arrow[dl,hook,"\underline{1}_\tau"'] \arrow[dr,hook,"\underline{0}_{\tau'}"] \& \\
I_\tau \arrow[dr,hook,"i_{\tau,\tau'}"'] \& \text{\unsure{cocart}} \& I_{\tau'} \arrow[dl,hook,"j_{\tau,\tau'}"] \\
\& I_{\tau\tau'} \& {}
\end{tikzcd}},
où toutes les flèches sont des foncteurs mono et pleinement fidèles. Que \unsure{le} diagramme \inserted{(évidemment commutatif)} soit bien cocartésien \inserted{dans $(\mathrm{Cat})$} revient, par définition, à dire que pour \deleted{\unsure{un}} toute $X$ dans $\mathrm{Cat}$, le diagramme transposé
\diagram{\begin{tikzcd}[ampersand replacement=\&,column sep=small,row sep=small]
\& \deleted{\unsure{\mathrm{Ch}_e|\mathrm{Ch}}}(X) \& \\
\mathrm{Ch}_\tau(X) \arrow[ur,phantom,"\deleted{\swarrow}"{description}] \arrow[ur,"\beta_\tau"{description},"\text{(but)}"'] \& \& \mathrm{Ch}_{\tau'}(X) \arrow[ul,phantom,"\deleted{\searrow}"{description}] \arrow[ul,"\leftarrow\!\sigma_{\tau'}\;\text{(source)}"']\\
\& \mathrm{Ch}_{\tau\tau'}(\unsure{X|\mathcal X}) \arrow[ul,phantom,"\deleted{\searrow}"{description}] \arrow[ur,phantom,"\deleted{\swarrow}"{description}] \arrow[ul] \arrow[ur] \&
\end{tikzcd}}
